\subsection{Conductors, collision singularities, and actual failure orbits}
Normalization alone does not describe the gluing of its branches.
Theorems~\ref{thm:split-conductor-v152} and
\ref{thm:binary-pinch-v152} determine two successive local models.
On the coprime two-divisor stratum the completed image is
\[
       \C[[z,x_1,\ldots,x_c,y_1,\ldots,y_c]]/(x_i y_j).
\]
Its conductor is $(x,y)$, its intrinsic singular ideal is $(x,y)^c$,
and its depth is $\dim(z)+1$.  When two common binary roots collide,
the model is instead
\[
 \C[[u,\Delta]]\oplus(u_1,\ldots,u_m)t
          \ \subset\ \C[[u,t]],\qquad \Delta=t^2,
\]
up to a specified smooth factor.  Its conductor has a ramified
double cover in the normalization.  These formulas determine the
Cohen--Macaulay cases and distinguish the singular scheme from its
reduced support.  Proposition~\ref{prop:degree-intersections-v152}
organizes the reduced intersections of different divisor-degree
images by the common part of two chosen divisors.

The boundary also has a direct role in the pencil inverse.
Theorem~\ref{thm:universal-boundary-v152} constructs a common point
in every intrinsic first-relation orbit closure of a pencil failure
algebra.  A scalar-matrix contraction of the flag-pencil relation
has gcd $a^{D-q_n}$, where $q_3=3$ and $q_n=4$ for $n\geq4$.
The selected degree-$n(n-1)$ divisor is unique but ramified; its
normalization fibre has tangent dimension
$\binom{n^2+3n-5-q_n}{3n-4-q_n}-1$.
Proposition~\ref{prop:pure-power-fibre-v152} gives its entire Artin
algebra by reciprocal-polynomial equations.  This specializes
\emph{actual} failure algebras through flat families; it does not
merely place an unrelated common-factor example beside the inverse
theorem.  The exact order of the flag coefficient filtration, proved
in Lemma~\ref{lem:flag-jets-v152}, is the link between the pencil
representation and the divisor collision.

\subsection{Factorization loci and the geometric comparison}
The normalization mechanism has classical precedents.
Chipalkatti's coincident-root locus parametrizes one binary form
with a prescribed root partition.  His Proposition~5.1, Theorem~5.4,
and Corollary~5.8 distinguish multiple normalization points from
failure of tangent injectivity and describe the singular support
\cite{Chipalkatti2003}.  Kurmann gives local incidence equations,
proves their scheme-theoretic equality in Theorem~1.5, and expresses
smoothness through splittings of partitions in Proposition~2.1
\cite{Kurmann2012}.  In particular, uniqueness plus tangent
injectivity is not a new general principle of the present paper.

Here the datum is an $r$-plane of forms in $e$ variables, not one
binary form; the selected divisor has prescribed degree, not a
root partition.  The residual datum is a Grassmannian subspace,
and for $e>2$ divisor degrees need not fill an interval.  These
changes produce high-codimension conductor strata and non-Cohen--Macaulay
branch gluings.  Theorems~\ref{thm:split-conductor-v152} and
\ref{thm:binary-pinch-v152} compute the image rings, not merely
the incidence graph or singular support.  Ferrand's conductor-square
formalism \cite{Ferrand2003} is the standard algebraic framework;
the particular conductor ideals, singular Fitting ideals, and depths
are calculated here for these Grassmannian strata.

There is a second, closer binary comparison.  Hu--Lin--Shao
\cite[\S2.2, Corollary~3.6, and Theorem~1.1]{HuLinShao2007}
study tuples of binary forms, define common-factor strata using
resultant homomorphisms, and resolve the map-space boundary by
successive blow-ups.  On the open set of linearly independent
$r$-tuples, passing to their span is the frame quotient onto
$\Gr(r,S_D)$; the common-divisor condition is invariant under this
quotient.  Thus our binary loci genuinely meet that theory rather
than forming an unrelated example.  We do not identify its
resultant scheme structure with our image structure without a
scheme comparison, nor claim a new general resolution.  Our
multivariate local rings and the pencil-orbit boundary are the
specific additional conclusions.

Rigidification is used as infrastructure, not as an independent
headline innovation.  The appropriate primary general result is
Abramovich--Olsson--Vistoli, Appendix~A, Theorem~A.1
\cite{AOV2008}, for compatible flat finitely presented normal
inertia subgroups, without a centrality assumption.  Our argument
identifies the particular substitution subgroup and twisted
right-factor subgroup and retains their actual group laws.  The
boundary statement is on the maximal-lower-Hilbert-function
first-relation stratum throughout.

The comparison with regular-pencil Segre strata is made in
Remark~\ref{rem:segre-boundary-v152}.  The common closed flag orbit
is not presented as a new classification of congruence orbits.
The new pencil-side conclusion is its explicit ramified boundary
under intrinsic cotangent changes, together with the complete
normalization-fibre algebra.  The separate documentary status of
Ballico's 1993 paper, discussed below, remains unchanged: absent
its full theorem text, no nonanticipation conclusion is drawn.
